\documentclass[11pt]{amsart}
\usepackage[margin=1.05in]{geometry}
\usepackage{amssymb,amsmath,booktabs,longtable,enumitem}
\usepackage{hyperref}

\newtheorem{theorem}{Theorem}
\newtheorem{lemma}[theorem]{Lemma}
\newtheorem{corollary}[theorem]{Corollary}
\theoremstyle{definition}
\newtheorem{definition}[theorem]{Definition}
\theoremstyle{remark}
\newtheorem{remark}[theorem]{Remark}

\newcommand{\cyc}{\mathcal{C}}
\newcommand{\Ext}{\operatorname{ext}}

\title[The short-path-and-cycles lemma]{The short-path-and-cycles lemma attributed to\\ Chu--Fan--Zhou: status, a proof of the common-vertex form, and its uses}
\date{26 November 2026}

\begin{document}

\begin{abstract}
The manuscripts of the program cite ``Chu, Fan and Zhou, Lemma~2.4'' in the form: a path $P$ with
$|V(P)|\le 4$ meeting each of $q$ edge-disjoint cycles, with $|V(P)|+q\le 7$, gives at most $q+1$
paths. We could not obtain a primary statement of that lemma. The only published statement of
this kind is Chu--Wang's own Lemma~2.2, which has the bound $|V(P)|+q\le 6$. We prove the bound
$|V(P)|+q\le 7$ under the hypothesis that all the cycles pass through one vertex of $P$
(Theorem~\ref{thm:main}). Every use with $|V(P)|+q=7$ in the manuscripts satisfies this
hypothesis. The proof is a short exchange argument. Its only inputs are the published path--cycle
lemma of Chu--Fan--Liu (at most five contacts) and a six-contact path--cycle lemma, which we verify
by a check of $1257$ compressed configurations. We do not prove the general form, in which the
cycles only meet $P$. We found no counterexample to it among $8{,}278{,}764$ configurations on at
most eight vertices.
\end{abstract}

\maketitle

\section{Summary}\label{sec:summary}

\begin{enumerate}[leftmargin=*]
\item \emph{Primary source.} We found no accessible copy of Y.~Chu, G.~Fan, C.~Zhou, ``Gallai's
conjecture and the path number of odd semi-cliques'', Discrete Math.\ 349(2) (2026) 114725
(DOI 10.1016/j.disc.2025.114725). The publisher's page returned 403. Unpaywall lists the paper
as closed. The arXiv API has no version under these authors or under the terms ``semi-clique''
and ``odd semi-clique''. Semantic Scholar and Crossref give only the metadata. The statement of
its Lemma~2.4 therefore remains \textbf{unverified}.

\item \emph{Chu--Wang Lemma~2.2 is not a quotation.} In arXiv~2510.12806v1 (source checked),
Chu and Wang introduce their Lemma~2.2 with ``We obtain the following lemma'' and prove it
themselves. Its hypothesis is $|V(P)|+|\cyc|\le 6$. They prove the case of six cycles through a
vertex separately, as their Lemma~2.3. Their reference [CFZ] is the 2022 triangle-free paper
(Discrete Math.\ 345, 112866), and they cite it only for its Theorem~3.1. So no verifiable source
states the $\le 7$ form, and no verifiable source states any form of Lemma~2.4 of the 2026 paper.

\item \emph{Every use with sum $7$ has a common vertex.} In every application in the manuscripts with
$|V(P)|+q=7$, all the cycles pass through a common vertex of $P$ (Section~\ref{sec:uses}).

\item \emph{What we prove.} Theorem~\ref{thm:main} is the common-vertex form with sum
$\le 7$. With $|V(P)|=1$ it gives Chu--Wang's Lemma~2.3 (six cycles through a vertex), with a
different proof.

\item \emph{What remains open.} The general form, in which each cycle only meets $P$, is not
proved. It holds in every configuration of the exhaustive checks in
Appendix~\ref{app:comp}. Those checks were run before the scope was narrowed and are not used in
any proof.
\end{enumerate}

\section{Conventions}

All graphs are finite and simple. A \emph{path} is a sequence of distinct vertices
$(v_1,\dots,v_k)$ with $k\ge 1$ and consecutive vertices adjacent. We allow $k=1$ only for the
path $P$ in the hypotheses. A member of a decomposition has at least one edge. Let $H$ be a graph
whose edge set is the disjoint union of the edge sets of some paths and cycles. We say that $H$
\emph{decomposes into $m$ paths} if $E(H)$ is the disjoint union of the edge sets of at most $m$
paths. For a cycle $C$ and $v\in V(C)$, $N_C(v)$ is the set of the two neighbours of $v$ on $C$.
For $t\in N_C(v)$, write $C-vt$ for the path obtained from $C$ by deleting the edge $vt$. It runs
from $v$ to $t$ and contains all of $V(C)$. This operation is called \emph{opening $C$ at $vt$}.

A \emph{configuration} $(P;C_1,\dots,C_q)$ consists of a path $P$ and cycles $C_i$ whose edge sets
are pairwise disjoint and disjoint from $E(P)$. Its \emph{union} is the graph formed by all these
edges, together with $V(P)$. In a configuration, two distinct edges at a vertex $v$ have distinct
other ends. In particular the sets $N_{C_i}(v)$ are pairwise disjoint and contain no
$P$-neighbour of $v$.

\section{Statements}

\begin{lemma}[path and cycle]\label{lem:pc}
Let $Q$ be a path and $C$ a cycle with $E(Q)\cap E(C)=\emptyset$ and $V(Q)\cap V(C)\ne\emptyset$.
Suppose $|V(Q)|\le 6$, and that $|V(Q)\cap V(C)|\le 4$ or $V(C)\not\subseteq V(Q)$. Then
$Q\cup C$ decomposes into two paths.
\end{lemma}

\begin{lemma}[Chu--Wang {\cite[Lemma 2.2]{CW}}]\label{lem:cw}
Let $(P;C_1,\dots,C_q)$ be a configuration in which every $C_i$ meets $V(P)$. If $|V(P)|\le 4$ and
$|V(P)|+q\le 6$, then the union decomposes into $q+1$ paths.
\end{lemma}

\begin{theorem}[short path and common cycles]\label{thm:main}
Let $(P;C_1,\dots,C_q)$ be a configuration and $x\in V(P)$ a vertex lying on every $C_i$. If
$|V(P)|\le 4$ and $|V(P)|+q\le 7$, then the union decomposes into $q+1$ paths.
\end{theorem}

\begin{corollary}\label{cor:ends}
In the setting of Theorem~\ref{thm:main}, suppose $q\ge 1$ or $P$ has an edge. If
$d=\deg_{H}(x)$ in the union $H$ satisfies $d\ge 2q+1$, then every decomposition into at most $q+1$
paths has exactly $q+1$ paths, all of them visit $x$, and exactly $2(q+1)-d$ of them end at $x$,
counted with multiplicity.
\end{corollary}

\begin{proof}
Each path contributes at most $2$ to $\deg_H(x)$, and it contributes $2$ exactly when it
traverses $x$. Every cycle passes through $x$, so $d=2q+\deg_P(x)\le 2q+2$. Count incidences at
$x$.
\end{proof}

\begin{corollary}[Chu--Wang {\cite[Lemma 2.3]{CW}}]\label{cor:six}
At most six edge-disjoint cycles through a common vertex decompose into one more path than their
number.
\end{corollary}

\begin{proof}
This is Theorem~\ref{thm:main} with $P=(x)$.
\end{proof}

\begin{remark}[the necessity of $|V(P)|\le 4$]
The condition $|V(P)|\le 4$ cannot be dropped, even with a common vertex. Take $P=(u_1,\dots,u_5)$
and the cycle $(u_1,u_3,u_5,u_2,u_4)$. This is the exceptional graph of Chu--Fan--Liu, which has
$9$ edges, while two paths on $5$ vertices have at most $8$.
\end{remark}

\subsection*{A Lean formulation}
With the definitions of \texttt{GallaiLean/Basic.lean} (\texttt{GPath}, \texttt{PathDecomp},
\texttt{pathNumber}), cycles through $x$ can be taken as closed walks based at $x$. Every cycle
through $x$ can be rotated to such a walk (\texttt{Walk.rotate}, \texttt{IsCycle.rotate}). The
needed statement is then:
\begin{verbatim}
theorem shortPath_commonCycles {V : Type*} [Fintype V] [DecidableEq V]
    (H : SimpleGraph V) {a b x : V} (P : H.Walk a b) (hP : P.IsPath)
    (hx : x ∈ P.support) {q : ℕ} (C : Fin q → H.Walk x x)
    (hC : ∀ i, (C i).IsCycle)
    (hnd : (P.edges ++ (List.ofFn fun i => (C i).edges).flatten).Nodup)
    (hcov : ∀ e, e ∈ H.edgeSet ↔
              e ∈ P.edges ++ (List.ofFn fun i => (C i).edges).flatten)
    (h4 : P.support.length ≤ 4) (h7 : P.support.length + q ≤ 7) :
    pathNumber H ≤ q + 1
\end{verbatim}
Following the treatment of \texttt{LovaszStatement}, the path--cycle input
(Lemma~\ref{lem:pc}) can enter as an explicit hypothesis:
\begin{verbatim}
def PathCycleStatement : Prop :=
  ∀ (W : Type u) [Fintype W] [DecidableEq W] (H : SimpleGraph W) {a b c : W}
    (Q : H.Walk a b), Q.IsPath → ∀ (C : H.Walk c c), C.IsCycle →
    (Q.edges ++ C.edges).Nodup →
    (∀ e, e ∈ H.edgeSet ↔ e ∈ Q.edges ++ C.edges) →
    (∃ v, v ∈ Q.support ∧ v ∈ C.support) → Q.support.length ≤ 6 →
    ((Q.support.toFinset ∩ C.support.toFinset).card ≤ 4 ∨
       ∃ v ∈ C.support, v ∉ Q.support) →
    pathNumber H ≤ 2
\end{verbatim}
Lemma~\ref{lem:cw} and Theorem~\ref{thm:main} then use only surgery on explicit paths. Mathlib's
characterisation \texttt{Walk.cons\_isCycle\_iff}, which says that \texttt{(cons h p).IsCycle} holds
iff \texttt{p.IsPath} and \texttt{s(u,v)} is not in \texttt{p.edges}, is exactly the opening
operation. The library's \texttt{PathDecomp.ofPerm} assembles a decomposition from explicit paths.
Decompositions of subgraphs transfer by \texttt{Walk.mapLe}. The six-contact case of
Lemma~\ref{lem:pc} is a finite statement about compressed configurations (Section~\ref{sec:pc}),
so it can be proved by \texttt{decide} once the compression step is formalized.

\section{Proof of Lemma~\ref{lem:pc}}\label{sec:pc}

Let $c=|V(Q)\cap V(C)|$.

\emph{Case $c\le 5$.} This is the lemma of Chu, Fan and Liu \cite[Lemma 2.1]{CFL}, quoted
verbatim as Lemma~2.1 of \cite{CW}. A connected union of an edge-disjoint path and cycle with at
most five common vertices decomposes into two paths, unless it is the exceptional graph. In the
exceptional graph $C$ is a $5$-cycle all of whose vertices lie on the path, as five consecutive
vertices, with a pendant subpath at each end of that segment. Our hypothesis excludes this:
either $c\le 4$, or $C$ has a vertex outside $Q$.

\emph{Case $c=6$.} Then $|V(Q)|=6$, $V(Q)\subseteq V(C)$, and $C$ has a vertex outside $Q$. The
six vertices of $Q$ cut $C$ into six arcs. Each arc is a single edge (\emph{direct}) or has
interior vertices, which lie outside $Q$ and have degree $2$ in $Q\cup C$. A direct arc is not an
edge of $Q$, by simplicity and edge-disjointness. At least one arc has interior vertices. Replace
every arc that has interior vertices by an arc with exactly one interior vertex, giving $C^\circ$.

A decomposition of $Q\cup C^\circ$ into two paths lifts to one of $Q\cup C$. A path through an
interior vertex $w$ traverses the whole arc. If a path ends at $w$, the other edge at $w$ lies on
the other path, which also ends at $w$. In that case give the first path the whole arc except its
last edge, and the second path that last edge. Both paths then end at the last interior vertex.
Interior vertices are private, so simplicity is preserved.

The compressed configurations are finite. Up to reflection there are $60$ cyclic orders of the
six labelled vertices of $Q=(0,\dots,5)$ on $C$. For each order we take every set of subdivided
arcs that contains all arcs parallel to an edge of $Q$ and is nonempty. This gives
$1257$ graphs, and each splits into two paths (Appendix~\ref{app:comp},
\texttt{path\_cycle\_check.py}). This agrees with the six-contact result in Section~6 of the
marked-reduction manuscript.

The same script checks every case of the lemma with $|V(Q)|\le 6$ independently of
\cite{CFL}. It enumerates the subset of contacts, their cyclic order, the subdivided arcs and the
off-cycle vertices of $Q$. The only configurations without a two-path partition have
$V(C)\subseteq V(Q)$: the pentagram (with or without a pendant vertex) at five contacts, and five
Hamiltonian $6$-cycles at six contacts. \qed

\section{Proof of Lemma~\ref{lem:cw}}\label{sec:cw}

We follow \cite{CW} and use Lemma~\ref{lem:pc} in place of \cite[Lemma 2.1]{CFL}. We need one
fact.

\emph{Fact.} For a path $(u_1,\dots,u_5)$, the only $5$-cycle on $\{u_1,\dots,u_5\}$ that is
edge-disjoint from it is $(u_1,u_3,u_5,u_2,u_4)$. Indeed, in $K_5$ minus the path, $u_2$, $u_3$
and $u_4$ have degree $2$. Their edges $u_2u_4$, $u_2u_5$, $u_3u_1$, $u_3u_5$, $u_4u_1$ are therefore
forced, and they already form the cycle.

We use induction on $q$. The cases $q=0$ and $q=1$ hold trivially and by Lemma~\ref{lem:pc} (at
most four contacts). Let $q\ge 2$ and $P=(v_1,\dots,v_k)$. Let $i$ be least such that $v_i$ lies on
a cycle, say $C_1$, and let $N_{C_1}(v_i)=\{x_1,x_2\}$.

Suppose first that $x_1\notin V(P)$. Then $P_1=(C_1-v_ix_1)+(v_i,v_{i-1},\dots,v_1)$ is a path,
because no $v_j$ with $j<i$ lies on $C_1$. Put $P'=(x_1,v_i,\dots,v_k)$. Every other cycle meets
$P$ only in $\{v_i,\dots,v_k\}$, so it meets $P'$. Also $|V(P')|+q-1=k-i+1+q\le 6$. If
$|V(P')|\le 4$, induction finishes. Otherwise $|V(P')|=5$, so $i=1$, $k=4$ and $q=2$.

If $x_1,x_2\in V(P)$, then $\{x_1,x_2\}\subseteq\{v_{i+2},\dots,v_k\}$. This forces $i=1$, $k=4$,
$\{x_1,x_2\}=\{v_3,v_4\}$ and $q=2$.

In both remaining cases write $\cyc=\{C_1,C_2\}$.
\begin{itemize}[leftmargin=*]
\item Suppose $|V(P')|=5$ with $P'=(x_1,v_1,v_2,v_3,v_4)$. If Lemma~\ref{lem:pc} does not apply to
$P'$ and $C_2$, then $C_2$ is a $5$-cycle on $V(P')$. By the Fact,
$C_2=(x_1,v_2,v_4,v_1,v_3)$, so $v_1v_3,v_1v_4\in E(C_2)$. Hence $x_2\notin V(P)\cup V(C_2)$.
Replace $x_1$ by $x_2$. Now $C_2$ has the vertex $x_1$ outside the new $P'$, and
Lemma~\ref{lem:pc} applies.
\item Suppose $\{x_1,x_2\}=\{v_3,v_4\}$. Let $x$ be the neighbour of $v_3$ on $C_1$ other than
$v_1$. Then $x\notin V(P)$. Put $P''=(x,v_3,v_2,v_1,v_4)$. The edges $xv_3$ and $v_1v_3$ lie on
$C_1$, and $C_1-\{xv_3,v_1v_4\}$ consists of the edge $v_1v_3$ and a path from $x$ to $v_4$ that
avoids $v_1$ and $v_3$. So $P_4=(C_1-xv_3-v_1v_4)+v_3v_4$ is the path
$(v_1,v_3,v_4,\dots,x)$. If Lemma~\ref{lem:pc} failed for $P''$ and $C_2$, then by the Fact
$C_2=(x,v_2,v_4,v_3,v_1)$. This cycle contains the edge $v_3v_4$ of $P$, which is impossible.
\end{itemize}
In each case we obtain $q+1=3$ paths. \qed

\section{Three operations}

Throughout, $(P;C_1,\dots,C_q)$ is a configuration and $x\in V(P)$ lies on every cycle.

\begin{lemma}[first contact]\label{lem:first}
Let $|V(P)|=k\le 4$, $P=(p_1,\dots,p_k)$, and suppose some cycle $C$ misses $p_1$. Let $p_j$ be
the first vertex of $P$ on $C$. Then $C$ has a neighbour $t$ of $p_j$ with
$t\notin V(P)$. The sequence $(C-p_jt)+(p_j,p_{j-1},\dots,p_1)$ is a path, and so is
$P'=(t,p_j,\dots,p_k)$. Moreover $P'$ contains $x$ and $|V(P')|\le k$.
\end{lemma}

\begin{proof}
Since $x\in V(C)$, we have $j\le$ the position of $x$, so $x\in V(P')$, and $j\ge 2$. The vertices
$p_1,\dots,p_{j-1}$ are not on $C$, and $p_{j+1}$ is a $P$-neighbour of $p_j$. So
$N_C(p_j)\cap V(P)\subseteq\{p_{j+2},\dots,p_k\}$, a set of at most $k-j-1\le 1$ vertex.
\end{proof}

\begin{lemma}[end opening]\label{lem:end}
Let $|V(P)|=k\le 3$ with $P=(p_1,\dots,p_k)$. Then every cycle $C$ through $p_1$ has a neighbour
$t\notin V(P)$ of $p_1$. Moreover $C-p_1t$ and $(t,p_1,\dots,p_k)$ are paths.
\end{lemma}

\begin{proof}
$N_C(p_1)\cap V(P)\subseteq V(P)\setminus\{p_1,p_2\}$, which has at most one element.
\end{proof}

\begin{lemma}[exchange]\label{lem:exch}
Let $v$ lie on cycles $C,D$ of a configuration.
\begin{enumerate}[label=(\alph*)]
\item If $u\in N_C(v)\setminus V(D)$ and $w\in N_D(v)\setminus V(C)$, then $(C-vu)+vw$ and
$(D-vw)+vu$ are two paths whose edge sets partition $E(C)\cup E(D)$.
\item Let $v$ lie on $C_1,C_2,C_3$ with $u_i\in N_{C_i}(v)\setminus V(C_{i+1})$ (indices mod $3$).
Then $(C_{i+1}-vu_{i+1})+vu_i$, $i=1,2,3$, are three paths partitioning
$E(C_1)\cup E(C_2)\cup E(C_3)$.
\end{enumerate}
\end{lemma}

\begin{proof}
$C-vu$ is a path from $v$ to $u$, and appending the edge $vw$ at $v$ keeps it simple because
$w\notin V(C)$. Part~(b) is the same argument applied three times.
\end{proof}

\section{Proof of Theorem~\ref{thm:main}}\label{sec:proof}

Let $k=|V(P)|$. If $k+q\le 6$, Lemma~\ref{lem:cw} applies. So let $k+q=7$, and argue by
descending induction on $k$, from $k=4$ to $k=1$.

\subsection*{Reduction of $k\le 3$ to $k+1$}
Let $k\le 3$, so $q\ge 4$. Suppose some cycle misses an end of $P$, say $p_1$. Apply
Lemma~\ref{lem:first}. The remaining $q-1$ cycles pass through $x\in V(P')$, and
$|V(P')|+q-1\le 6$. Lemma~\ref{lem:cw} gives $q$ paths, and the opened cycle is one more. Note
that $k=1$ never reaches this case.

Otherwise all cycles contain $p_1$. Pick any cycle $C$ and apply Lemma~\ref{lem:end}. The path
$P'=(t,p_1,\dots,p_k)$ has $k+1\le 4$ vertices and contains $x$. The other $q-1$ cycles pass
through $x$, and $(k+1)+(q-1)=7$. By induction they and $P'$ give $q$ paths, and $C-p_1t$ is the
last one.

\subsection*{The case $k=4$, $q=3$}
Write $P=(p_1,p_2,p_3,p_4)$, let $H$ be the union, and call the cycles $C_1,C_2,C_3$. Each of
the edges $p_1p_3$, $p_1p_4$, $p_2p_4$ lies on at most one cycle. For every cycle $C$,
$$N_C(p_1)\cap V(P)\subseteq\{p_3,p_4\},\qquad N_C(p_4)\cap V(P)\subseteq\{p_1,p_2\}.$$

\emph{Step 1.} Suppose a cycle misses $p_1$ or $p_4$, say $p_1$ (otherwise reverse $P$). By
Lemma~\ref{lem:first}, the other two cycles and $P'$ satisfy Lemma~\ref{lem:cw}, since
$|V(P')|\le 4$ and $|V(P')|+2\le 6$. They give $3$ paths, so $4$ in all. \emph{From now on every
cycle contains $p_1$ and $p_4$.} The vertex $x$ is not used again.

\emph{Step 2 (exchanges at $p_1$).} Write $i\to j$ if $N_{C_i}(p_1)\not\subseteq V(C_j)$. Suppose
$i\to j$ and $j\to i$. Then Lemma~\ref{lem:exch}(a) gives two paths on $C_i\cup C_j$.
Lemma~\ref{lem:pc} gives two paths on $P\cup C_k$, because they share $p_1$ and have at most
$|V(P)|=4$ common vertices. This makes $4$. If $1\to2\to3\to1$ or $1\to3\to2\to1$,
Lemma~\ref{lem:exch}(b) gives three paths, and $P$ is the fourth.

Otherwise the relation $\to$ is an acyclic digraph on three vertices, so some cycle $S$ has no
incoming arc. Then $N_C(p_1)\subseteq V(S)$ for every cycle $C$. The six neighbours of $p_1$ on the
cycles are distinct, so $|V(S)|\ge 7$.

\emph{Step 3 (two-end extension).}
\begin{lemma}\label{lem:two}
Let $\{F,X,Y\}=\{C_1,C_2,C_3\}$, and let $t\in N_X(p_1)\setminus V(P)$ and
$y\in N_Y(p_4)\setminus(V(P)\cup\{t\})$. If $V(F)\not\subseteq\{t,p_1,p_2,p_3,p_4,y\}$, then
$H$ decomposes into four paths.
\end{lemma}
\begin{proof}
$X-p_1t$ and $Y-p_4y$ are paths, and so is $R=(t,p_1,p_2,p_3,p_4,y)$. The cycle $F$ contains $p_1$,
$|V(R)|=6$, and $F$ has a vertex outside $R$. Lemma~\ref{lem:pc} gives two paths on $F\cup R$. The
four edge sets partition $E(H)$.
\end{proof}

\emph{Step 4.} Let $A,B$ be the two cycles other than $S$. Put
$$\alpha=N_A(p_1)\setminus V(P),\quad \beta=N_B(p_1)\setminus V(P),\quad
\alpha'=N_A(p_4)\setminus V(P),\quad \beta'=N_B(p_4)\setminus V(P).$$
Take $F=S$ in Lemma~\ref{lem:two}. The condition on $F$ holds automatically, since $|V(S)|\ge 7$.
\begin{itemize}[leftmargin=*]
\item \emph{Option 1} takes $(X,Y)=(A,B)$. It works unless $\alpha=\emptyset$, or $\beta'=\emptyset$,
or $\alpha=\beta'$ is a single vertex.
\item \emph{Option 2} takes $(X,Y)=(B,A)$. It works unless $\beta=\emptyset$, or $\alpha'=\emptyset$, or
$\beta=\alpha'$ is a single vertex.
\end{itemize}
Note that $\alpha=\emptyset$ means $N_A(p_1)=\{p_3,p_4\}$. In that case $B$ and $S$ use neither
$p_1p_3$ nor $p_1p_4$, so $|\beta|=2$. The same holds at $p_4$ and for $B$.

\emph{Claim.} If both options fail, then up to interchanging $A$ and $B$ one of the following
holds.
\begin{itemize}[leftmargin=*]
\item[(F1)] $N_A(p_1)=\{p_3,p_4\}$ and $N_A(p_4)=\{p_1,p_2\}$.
\item[(F3)] $p_1p_4\in E(A)$, and for distinct vertices $v,v'\notin V(P)$,
$$N_A(p_1)=\{v,p_4\},\quad N_A(p_4)=\{p_1,v'\},\quad N_B(p_1)=\{v',p_3\},\quad N_B(p_4)=\{v,p_2\}.$$
\end{itemize}

\begin{proof}[Proof of the claim]
Suppose $\alpha=\emptyset$. Then $|\beta|=2$, so Option~2 can fail only through
$\alpha'=\emptyset$, which gives (F1). The case $\beta'=\emptyset$ is symmetric: $|\alpha'|=2$, so
$\beta=\emptyset$, which is (F1) with $A$ and $B$ interchanged.

Otherwise $\alpha$ and $\beta'$ are nonempty, so $\alpha=\beta'=\{v\}$. Then $A$ uses exactly one
of $p_1p_3$, $p_1p_4$, and $B$ uses exactly one of $p_4p_1$, $p_4p_2$. Hence $\beta\ne\emptyset$
and $\alpha'\ne\emptyset$, so Option~2 fails through $\beta=\alpha'=\{v'\}$. Thus each of $A,B$
has exactly one $P$-neighbour at $p_1$ and one at $p_4$. These are distinct edges, so
$\{p_3,p_4\}$ is split between $A$ and $B$ at $p_1$, and $\{p_1,p_2\}$ at $p_4$. The edge
$p_1p_4$ lies on $A$ or on $B$. Interchanging $(A,v)$ with $(B,v')$ if necessary, it lies on $A$,
which is (F3). Finally $v\ne v'$, since $p_1v$ and $p_1v'$ are distinct edges.
\end{proof}

\emph{Resolving (F3).} The edges $p_1p_3$ and $p_1p_4$ lie on $B$ and $A$. So both neighbours of
$p_1$ on $S$ lie outside $P$. Pick $t\in N_S(p_1)$ and apply Lemma~\ref{lem:two} with $F=A$,
$X=S$, $Y=B$ and $y=v$. Here $t\ne v$, since they are the ends of distinct edges at $p_1$, and
$v\notin V(P)$. The vertex $v'\in N_A(p_4)\subseteq V(A)$ lies outside
$\{t,p_1,\dots,p_4,v\}$: it is not in $V(P)$, it differs from $v$, and it differs from $t$, since
$p_1v'\in E(B)$ and $p_1t\in E(S)$.

\emph{Resolving (F1).} $A$ contains the path $(p_3,p_1,p_4,p_2)$. The rest of $A$ is a path from $p_2$
to $p_3$ that is not the edge $p_2p_3$. Its interior $W^\circ\ne\emptyset$ lies outside $V(P)$,
and $V(A)=V(P)\cup W^\circ$. The cycles $B$ and $S$ use none of $p_1p_3$, $p_1p_4$, $p_2p_4$. Hence
$$N_B(p_1)=\{b_1,b_2\},\quad N_S(p_1)=\{s_1,s_2\},\quad N_B(p_4)=\{d_1,d_2\},\quad N_S(p_4)=\{e_1,e_2\},$$
all outside $V(P)$, with the four vertices $b_i,s_i$ distinct and the four vertices $d_i,e_i$
distinct. We apply Lemma~\ref{lem:two} with $F=A$, using either
$(X,Y)=(B,S)$ with $t\in\{b_1,b_2\}$, $y\in\{e_1,e_2\}\setminus\{t\}$, or $(X,Y)=(S,B)$ with
$t\in\{s_1,s_2\}$, $y\in\{d_1,d_2\}\setminus\{t\}$. It suffices that $W^\circ\not\subseteq\{t,y\}$.
\begin{itemize}[leftmargin=*]
\item If $|W^\circ|\ge 3$, any choice works.
\item Let $|W^\circ|=2$. Two admissible pairs $(b_i,e_j)$ that share a coordinate give different sets
$\{t,y\}$, so one of them differs from $W^\circ$. Otherwise at most two pairs are admissible. This
happens only when $\{b_1,b_2\}=\{e_1,e_2\}$, and then the only set is $\{b_1,b_2\}$. If that set is
$W^\circ$, use $(S,B)$ with $t=s_1\notin W^\circ$.
\item Let $W^\circ=\{w\}$. If $w\notin\{b_1,b_2\}$, then $\{e_1,e_2\}\setminus\{w\}\ne\emptyset$.
Take $y$ in it and $t\in\{b_1,b_2\}\setminus\{y\}$. If $w\in\{b_1,b_2\}$, use $(S,B)$: $t=s_i$ is
never $w$, and $\{d_1,d_2\}\setminus\{w\}$ contains a vertex $y\ne t$ for a suitable
$i\in\{1,2\}$.
\end{itemize}
This completes the case $k=4$, $q=3$, and with it the proof of Theorem~\ref{thm:main}. \qed

\begin{remark}
In the proof, $x$ is used only to make Step~1 and the reduction to $k+1$ valid. There it guarantees
that the shortened path still meets every remaining cycle. The general form, in which each cycle
only meets $P$, would follow from the same argument if Step~1 could be replaced. That is the
configuration in which some cycle misses an end of $P$ and the other cycles meet $P$ only before
that cycle's first contact. We have no proof of that case.
\end{remark}

\begin{remark}[relation to the internal results]
The marked-reduction manuscript contains Lemma~173 (a path $(A,B,C,V,Z)$ with four cycles through
$V$ gives five paths). Padding a path with private pendant edges, and adding dummy triangles
through the common vertex, removes edges only from path ends. So Lemma~173 also implies
Theorem~\ref{thm:main}, and more: it gives $|V(P)|+q\le 8$ for $|V(P)|\le 4$. Its proof, however,
rests on the computer-checked Theorem~159. Theorem~38 there (at most seven cycles through a
vertex give $j+1$ paths with two ends at the vertex) does not use Lemma~2.4.
\end{remark}

\section{Uses in the manuscripts and their status}\label{sec:uses}

Line numbers refer to the text extractions. Here ``(a)'' means covered by Lemma~\ref{lem:cw}
(sum $\le 6$), ``Thm~\ref{thm:main}'' means sum $7$ with a common vertex, now proved here, and
``(b)'' means covered by marked-reduction Theorem~38, which is independent of Lemma~2.4. In every
case the endpoint conclusions drawn in the manuscripts follow from degree counting at the common
vertex (Corollary~\ref{cor:ends}).

\begin{longtable}{p{3.1cm}p{5.1cm}p{1.3cm}p{4.3cm}}
\toprule
Place & Configuration & $(k,q)$ & Status\\
\midrule
\endhead
MR L.~21 ($\approx$2181), case 1 & $P=(a,b,x)$, $q\le 3$ cycles through $x$ & $(3,\le3)$ & (a)\\
MR L.~21, case 2 & $P=(a,b,x,c)$; $q\le 3$, or $q\le 2$ at $d=8$ & $(4,3)$ & Thm~\ref{thm:main} ($x$ common)\\
MR L.~21, case 3 & $P=(x,a,b,c)$, $q\le 3$ & $(4,3)$ & Thm~\ref{thm:main}\\
MR L.~21, rank remark ($\approx$2189), used at $\approx$2805 & case 1 with four cycles & $(3,4)$ & Thm~\ref{thm:main}\\
MR Prop.~27 ($\approx$2466) & $P=(a,b,x,z)$, $\le1$ cycle & $(4,\le1)$ & (a)\\
MR Prop.~33 ($\approx$2755) & planted 4-vertex path; 3 closed Hamilton rows of $U+x$ & $(4,3)$ & Thm~\ref{thm:main}: every row contains $V(P)$, so any $p_i$ is common\\
MR Cor.~35 ($\approx$2789) & reserved spoke $xz$; $q\le5$ closed rows through $x$ & $(2,\le5)$ & (a) for $q\le4$; Thm~\ref{thm:main} for $q=5$\\
MR Cor.~35 planting ($\approx$2840) & $P$ of order $L+b+1=4$ through $x$; $q\le 3$ rows through $x$ & $(4,3)$ & Thm~\ref{thm:main}\\
MR mixed pair ($\approx$2929) & $(a,b,x)$, $\le3$ cycles & $(3,\le3)$ & (a)\\
MR Thm.~69 ($\approx$4690) & spoke or $2$-edge path at $x$, $\le2$ cycles & $\le(3,2)$ & (a)\\
MR \S17.3 ($\approx$4749), L.~70, Thm.~72 ($\approx$4891, 4899) & stated as ``$\le7$''; used as $(3,\le2)$, $(4,\le1)$, $(2,\le3)$, $(3,\le3)$ & sums $\le6$ & (a). The stated general sentence is exactly Thm~\ref{thm:main} (cycles meet $P$ at $x$)\\
MR L.~108 ($\approx$6285) & edge $+\le2$ cycles; 3-edge path $+1$ cycle & $(2,2)$, $(4,1)$ & (a)\\
MR Thm.~177 ($\approx$9752) & $P=(b,a,v,x)$, $m\le3$ cycles through $v,x$ & $(4,3)$ & Thm~\ref{thm:main} ($v$ common)\\
MR $\approx$3185 (``Lemma 2.5'', $j\le7$) & cycles through a vertex & $(1,\le7)$ & not load-bearing; (b)\\
even-degrees 137 (``Lemma 2.5'', $m\le7$) & $m$ cycles through $x$ & $(1,\le7)$ & Cor.~\ref{cor:six} for $m\le6$; (b) for $m=7$\\
even-degrees 156--157 & edge $xy$ and $m-1\le5$ cycles through $x$ & $(2,\le5)$ & (a) for $m\le5$; Thm~\ref{thm:main} (also (b)) for $m=6$\\
separating-pairs 3934--3936 & route of $\le4$ vertices, $\le2$ cycles & sum $\le6$ & (a)\\
degree-ten-marked-root 60, 560 & 4-vertex path with $\le2$ cycles; $\le3$-vertex path with $\le2$ cycles & sum $\le6$ & (a)\\
four-exterior-neighbors 589, 597, 779; five-odd-neighbors 117, 449 & internal MR Lemmas 152/155, $\le2$ cycles & sum $\le6$ & not citations of Lemma 2.4\\
\bottomrule
\end{longtable}

\noindent\emph{Recommended edits.} In each sum-$7$ place, replace ``Chu, Fan and Zhou [8, Lemma
2.4]'' by Theorem~\ref{thm:main} of this note, and state the common-vertex hypothesis, which is
already true there. The citation of Lemma~2.4 elsewhere may be replaced by Lemma~\ref{lem:cw}
(Chu--Wang, Lemma~2.2). Where ``Lemma 2.5'' is cited for $m=7$, cite marked-reduction Theorem~38
instead.

\appendix
\section{Computations}\label{app:comp}

The scripts are in \texttt{scripts/} next to this file and the outputs in \texttt{data/}. The only
computation that the proofs depend on is the six-contact part of A.1, which takes under a second.
Everything else is a sanity check.

\begin{enumerate}[label=A.\arabic*, leftmargin=*]
\item \texttt{path\_cycle\_check.py} (with the exact oracle \texttt{decomp.py}) verifies
Lemma~\ref{lem:pc} for $|V(Q)|\le6$ on compressed configurations. Output:
\texttt{data/path\_cycle\_check.out}. At six contacts with an outside vertex there are $1257$
configurations and $0$ failures. All failures in the whole range have $V(C)\subseteq V(Q)$: the
pentagram at five contacts, and five Hamiltonian $6$-cycles at six.

\item \texttt{proof43\_check.py} and \texttt{proof43\_targeted.py} run the case analysis of
Section~\ref{sec:proof} as an algorithm on random configurations. These include triples of Hamilton
cycles of $K_9,K_{11},K_{13}$ (from a Walecki decomposition), configurations nested near $P$, and
instances built to reach (F1) and (F3). Each branch builds its paths explicitly, and the steps
that invoke Lemma~\ref{lem:pc} or Lemma~\ref{lem:cw} are checked with the exact oracle. No
assertion failed. Every branch, including (F1) with both orders and (F3), was reached.

\item \texttt{flower.c} is an exhaustive check of the \emph{general} form, in which each cycle only
meets $V(P)$. It takes all sets of $q$ edge-disjoint cycles of $K_n-E(P)$ with
$P=(0,\dots,k-1)$ and decides by exact search whether $q+1$ paths suffice. Results
(\texttt{data/exh\_*}):
\begin{itemize}
\item $(4,3)$: $n=7$ ($18{,}222$ sets) and $n=8$ ($8{,}278{,}764$ sets), with $0$ failures.
\item At $n=7$, $0$ failures for $(3,4)$ ($18{,}699$ sets), $(2,4)$ ($29{,}430$), $(5,2)$
($4{,}248$), $(5,3)$, $(4,4)$ and $(3,5)$.
\item The only failures found are for $(6,1)$ and $(6,2)$. They are the pentagram, Hamiltonian
$6$-cycles on $V(P)$, and pairs of Hamilton cycles of $K_7$, which the edge count rules out.
\end{itemize}
This run predates the narrowing of the task to conceptual arguments. It is supporting evidence
only.

\item \texttt{rflower.c} runs random tests with vertex pools of up to $12$ vertices, cycle lengths
up to $10$, and both the general and the common-vertex mode. No failure was found.
\end{enumerate}

\begin{thebibliography}{9}
\bibitem{CFL} Y.~Chu, G.~Fan, Q.~Liu, On Gallai's conjecture for graphs with maximum degree 6,
Discrete Math.\ 344 (2021) 112212, Lemma~2.1 (as quoted in \cite{CW}).
\bibitem{CFZ} Y.~Chu, G.~Fan, C.~Zhou, Gallai's conjecture and the path number of odd
semi-cliques, Discrete Math.\ 349(2) (2026) 114725. Not consulted: no copy was accessible.
\bibitem{CW} Y.~Chu, Y.~Wang, Path decompositions of Eulerian graphs, arXiv:2510.12806v1 (2025),
Lemmas~2.1--2.3.
\end{thebibliography}

\end{document}
